% concatenated from main.tex
\documentclass[11pt]{article}
\usepackage[a4paper,margin=1in]{geometry}
\usepackage[T1]{fontenc}
\usepackage{lmodern}
\usepackage{amsmath,amssymb,mathtools,amsthm}
\usepackage{graphicx}
\usepackage{listings}
\usepackage{microtype}
\usepackage{cite}
\setlength{\parindent}{1.15em}
\setlength{\parskip}{0pt}
\setlength{\abovedisplayskip}{6pt plus 1pt minus 1pt}
\setlength{\belowdisplayskip}{6pt plus 1pt minus 1pt}
\setlength{\abovedisplayshortskip}{4pt}
\setlength{\belowdisplayshortskip}{4pt}
\newtheorem{lemma}{Lemma}
\newtheorem{theorem}{Theorem}
\newtheorem{proposition}{Proposition}
\DeclareMathOperator{\sgn}{sgn}
\newcommand{\ii}{\mathrm i}
\newcommand{\dd}{\,\mathrm d}
\newcommand{\Rea}{\operatorname{Re}}
\newcommand{\Ima}{\operatorname{Im}}
\newcommand{\br}[3]{\left[#1\right]_{#2}^{#3}}
\lstset{
  basicstyle=\ttfamily\scriptsize,
  breaklines=true,
  columns=fullflexible,
  keepspaces=true,
  showstringspaces=false,
  literate={ü}{{\"u}}1 {é}{{\'e}}1 {±}{{$\pm$}}1,
  aboveskip=6pt,
  belowskip=6pt
}

\title{\vspace{-12mm}\textbf{Explicit Cartesian Coordinates for Three-Dimensional Clothoids}}
\author{Alexandru Ionu\textcommabelow{t}}
\date{}

\begin{document}
\maketitle
\vspace{-6mm}
\begin{abstract}
We combine the spherical-clothoid reconstruction theorem of Lucas and Ortega-Yag\"ues with an explicit spinorial parametrization of the spherical clothoid and a Kamp\'e de F\'eriet evaluation of Kummer-product integrals. This gives explicit Cartesian coordinates, in arclength, for every nonhelical three-dimensional clothoid. Kummer asymptotics then yield explicit asymptotic lines, including their gamma--digamma offsets and the angle between their axes.
\end{abstract}
\noindent\textbf{Keywords.} 3D clothoid; spherical clothoid; confluent hypergeometric function; Kamp\'e de F\'eriet function.

\section{Introduction}
A three-dimensional clothoid is a unit-speed space curve whose curvature and torsion are affine functions of arclength,
\begin{equation}
\kappa(s)=cs+d,\qquad \tau(s)=as+b. \tag{1}
\end{equation}
Throughout, \(\kappa(s)=cs+d\) denotes the signed coefficient in the smoothly continued Frenet equations; away from its zeros, the ordinary curvature is \(|cs+d|\).
Berry and Robbins earlier obtained an exact spinorial parametrization for the special case of constant curvature and linearly varying torsion, through its equivalence with the Landau--Zener problem \cite{BerryRobbins1993}.
Frego obtained closed Fresnel formulas for an analytically tractable subfamily and treated the remaining cases numerically \cite{Frego2022}. More recently, Lucas and Ortega-Yag\"ues proved that every nonhelical curve satisfying (1) can be reconstructed from a spherical clothoid, but their formula retains one vector quadrature of the spherical normal \cite[Theorem~4.4]{LucasOrtega2026}.

The Sabban equations admit a two-component spinor formulation \cite{Senyurt2015}, and the spherical clothoid together with its full Sabban frame has an explicit representation in confluent hypergeometric functions \cite{Ionut2022}. Substitution of that frame into the reconstruction theorem turns the formula itself into a single integral of spinor quadratics. The remaining Kummer-product integrals are then evaluated by the Kamp\'e de F\'eriet series studied by Jur\v{s}\.{e}nas \cite{Jursenas2014}.

\section{The spinorial integral form of Theorem 4.4}
Following \cite{LucasOrtega2026}, put
\begin{equation}
\Delta=ad-bc,\qquad \varepsilon=\sgn\Delta,\qquad \rho=\sqrt{a^2+c^2}, \tag{2}
\end{equation}
and assume \(\Delta\ne0\). The constants in their Theorem~4.4 are
\begin{equation}
 p=\frac{\varepsilon a\Delta}{\rho^3},\quad
 q=-\frac{\varepsilon c}{\rho},\quad
 \lambda=\frac{\varepsilon\Delta}{\rho},\quad
 m=\frac{\varepsilon\rho^3}{\Delta^2},\quad
 n=\frac{ab+cd}{\Delta}. \tag{3}
\end{equation}
If \(\delta\) is the unit-speed spherical clothoid with geodesic curvature \(\kappa_g(u)=mu+n\) and intrinsic spherical normal \(N_\delta\), then, up to an orientation-preserving Euclidean congruence,
\begin{equation}
 r(s)=p\bigl(\kappa_g(\lambda s)\delta(\lambda s)+N_\delta(\lambda s)\bigr)
      +q\int_{s_0}^{s}N_\delta(\lambda u)\dd u. \tag{4}
\end{equation}
Introduce
\begin{equation}
 \chi=-\frac{\Delta}{\rho^{3/2}},\qquad
 \xi(s)=\sqrt\rho\left(s+\frac{ab+cd}{\rho^2}\right),\qquad \xi_0=\xi(0), \tag{5}
\end{equation}
so that
\begin{equation}
 \kappa_g(\lambda s)=m\lambda s+n=-\frac{\xi(s)}{\chi}. \tag{6}
\end{equation}
Set
\begin{equation}
 \mu=\frac12+\frac{\ii\chi^2}{8},\qquad \zeta=\frac{\ii\xi^2}{2},\qquad
 U(\xi)={}_1F_1\!\left(\mu;\frac12;\zeta\right),\qquad
 V(\xi)={}_1F_1\!\left(\mu;\frac32;\zeta\right), \tag{7}
\end{equation}
\begin{equation}
 \psi_1(\xi)=e^{-\zeta/2}U(\xi),\qquad
 \psi_2(\xi)=\frac{\ii\chi\xi}{2}e^{-\zeta/2}V(\xi). \tag{8}
\end{equation}
For \(v=(v_1,v_2,v_3)\in\mathbb R^3\), write \(v^\perp:=v_1+\ii v_2\), while \(v_3\) denotes its third component. We use the Sabban ordering
\(\{\gamma,\mathbf t,\mathbf d\}\) and convention \(\mathbf d=\gamma\times\mathbf t\) of
\cite{Senyurt2015}. Equivalently, the spinor representing
\(\{\mathbf t,\mathbf d,\gamma\}\) in the convention of that paper is
\(e^{-\pi\ii/4}(\overline{\psi_1},-\psi_2)^T\). The complete frame is
\begin{align}
 \gamma^\perp&=-2\psi_1\psi_2,& \gamma_3&=|\psi_1|^2-|\psi_2|^2, \tag{9}\\
 \mathbf t^\perp&=\ii(\psi_1^2+\psi_2^2),& (\mathbf t)_3&=2\Ima(\overline{\psi_1}\psi_2), \notag\\
 \mathbf d^\perp&=-(\psi_1^2-\psi_2^2),& (\mathbf d)_3&=-2\Rea(\overline{\psi_1}\psi_2). \tag{10}
\end{align}
Choose the orientation so that the notation of \cite{LucasOrtega2026} is
\begin{equation}
 \delta(\lambda s)=\varepsilon\gamma(\xi(s)),\qquad
 T_\delta(\lambda s)=\mathbf t(\xi(s)),\qquad
 N_\delta(\lambda s)=\varepsilon\mathbf d(\xi(s)). \tag{11}
\end{equation}
In the variable \(\xi\), the pulled-back Sabban equations are
\begin{equation}
 \gamma'=-\chi\mathbf t,\qquad
 \mathbf t'=\chi\gamma+\xi\mathbf d,\qquad
 \mathbf d'=-\xi\mathbf t,
 \qquad\text{hence}\qquad (\xi\gamma-\chi\mathbf d)'=\gamma. \tag{12}
\end{equation}
Substituting (3)--(11) into (4), translating so that \(r(0)=0\), and using (12), gives the direct integral form
\begin{equation}
 \boxed{\displaystyle
 r(s)=\rho^{-3/2}\int_{\xi_0}^{\xi(s)}\bigl(a\,\gamma(\eta)-c\,\mathbf d(\eta)\bigr)\dd\eta.} \tag{13}
\end{equation}
Thus Theorem~4.4 is already an explicit spinorial quadrature. In Cartesian form,
\begin{align}
 x(s)+\ii y(s)
 &=\rho^{-3/2}\int_{\xi_0}^{\xi(s)}
 \Bigl[c(\psi_1^2-\psi_2^2)-2a\psi_1\psi_2\Bigr]\dd\eta, \tag{14}\\
 z(s)
 &=\rho^{-3/2}\int_{\xi_0}^{\xi(s)}
 \Bigl[2c\Rea(\overline{\psi_1}\psi_2)+a(|\psi_1|^2-|\psi_2|^2)\Bigr]\dd\eta, \tag{15}
\end{align}
where the spinors in the integrands are evaluated at \(\eta\). Equation (12) also shows that only the normal-field part requires a new primitive:
\begin{equation}
 r(s)=\rho^{-3/2}\br{cJ_\chi(\xi)+a\bigl(\xi\gamma(\xi)-\chi\mathbf d(\xi)\bigr)}{\xi_0}{\xi(s)},
 \qquad J_\chi'=-\mathbf d. \tag{16}
\end{equation}

\section{Evaluation of the normal-field primitive}
We use the convention
\begin{equation}
F_{1:1;1}^{1:1;1}\!\left[
\begin{matrix}A:B;C\\D:E;F\end{matrix};X,Y\right]
=\sum_{j,k\ge0}\frac{(A)_{j+k}(B)_j(C)_k}{(D)_{j+k}(E)_j(F)_k}\frac{X^jY^k}{j!k!}. \tag{17}
\end{equation}
Define
\begin{align}
\mathcal F_\chi(\xi)
&:=F_{1:1;1}^{1:1;1}\!\left[
\begin{matrix}\frac12:\mu;\frac12-\mu\\[1mm]\frac32:\frac12;\frac12\end{matrix};\zeta,-\zeta\right], \tag{18}\\
\mathcal G_\chi(\xi)
&:=F_{1:1;1}^{1:1;1}\!\left[
\begin{matrix}1:\mu;1-\mu\\[1mm]2:\frac32;\frac12\end{matrix};\zeta,-\zeta\right]. \tag{19}
\end{align}
\begin{lemma}
For real \(\xi\),
\begin{equation}
 (\xi\mathcal F_\chi)'=\psi_1^2,\qquad
 \left(\frac{\ii\chi\xi^2}{4}\mathcal G_\chi\right)'=\overline{\psi_1}\psi_2. \tag{20}
\end{equation}
Consequently,
\begin{equation}
 J_\chi^\perp=2\xi\mathcal F_\chi-\xi e^{-\zeta}UV,
 \qquad
 (J_\chi)_3=\Rea\left(\frac{\ii\chi\xi^2}{2}\mathcal G_\chi\right). \tag{21}
\end{equation}
\end{lemma}
\begin{proof}
The identities in (20) are the relevant specializations of the Kummer-product integral in \cite{Jursenas2014}; they also follow by multiplying the two \({}_1F_1\) series and integrating termwise. For \(N=j+k\),
\[
 \frac{(1/2)_N}{(3/2)_N}=\frac1{2N+1},\qquad
 \frac{(1)_N}{(2)_N}=\frac1{N+1},
\]
which cancel the derivatives of \(\xi^{2N+1}\) and \(\xi^{2N+2}/2\). Kummer's transformation, together with conjugation for real \(\xi\), gives the opposite-argument factors in (20). Finally,
\[
 (\psi_1\psi_2)'=\frac{\ii\chi}{2}(\psi_1^2+\psi_2^2)
\]
eliminates the third apparent primitive and yields (21).
\end{proof}

\section{Explicit Cartesian coordinates}
\begin{theorem}
Let (1) hold with \(a^2+c^2>0\) and \(\Delta\ne0\). With the notation (2), (5), (7), (18), and (19), the corresponding three-dimensional clothoid is, up to an orientation-preserving rigid motion,
\begin{align}
 x(s)+\ii y(s)
 &=\rho^{-3/2}\br{
 2c\xi\mathcal F_\chi+e^{-\zeta}\left\{
 a\chi U^2+\frac{a\chi^3\xi^2}{4}V^2-(c\xi+\ii a\chi\xi^2)UV
 \right\}}{\xi_0}{\xi(s)}, \tag{22}\\[1mm]
 z(s)
 &=\rho^{-3/2}\br{
 \Rea\left(\frac{\ii c\chi\xi^2}{2}\mathcal G_\chi\right)
 +a\xi\left\{2|U|^2-1-\chi^2\Ima(\overline{U}V)\right\}}{\xi_0}{\xi(s)}. \tag{23}
\end{align}
Every function in the brackets is evaluated at \(\xi\), and \(\br{f}{\xi_0}{\xi(s)}=f(\xi(s))-f(\xi_0)\).
\end{theorem}
\begin{proof}
Insert (9)--(10) and (21) into (16). For the third coordinate use \(|\psi_1|^2+|\psi_2|^2=1\) and
\(2\Rea(\overline{\psi_1}\psi_2)=-\chi\xi\Ima(\overline{U}V)\).
\end{proof}
The endpoint subtraction gives \(r(0)=0\); arbitrary initial position and orientation are obtained by a single rigid motion.

\section{Explicit asymptotic lines}
The integral representation (14)--(16), rather than the expanded Kamp\'e de F\'eriet expressions (22)--(23), is the convenient starting point for asymptotics. Define only the phase
\begin{equation}
 \Phi(\xi)=\frac{\xi^2}{4}+\frac{\chi^2}{4}\log|\xi|, \tag{24}
\end{equation}
and the two Kummer-branch coefficients
\begin{equation}
 A_\chi=
 \frac{\sqrt\pi\,e^{-\pi\chi^2/16}2^{-\ii\chi^2/8}}
 {\Gamma(\frac12+\frac{\ii\chi^2}{8})},\qquad
 B_\chi=
 \frac{\chi\sqrt\pi}{2\sqrt2}\,
 \frac{e^{-\pi\chi^2/16}2^{\ii\chi^2/8}e^{3\pi\ii/4}}
 {\Gamma(1-\frac{\ii\chi^2}{8})}. \tag{25}
\end{equation}

\begin{lemma}[Spinor asymptotics]
As \(\xi\to\pm\infty\), the spinors possess the completely explicit Poincar\'e expansions
\begin{align}
 \psi_1(\xi)\sim{}&A_\chi e^{\ii\Phi(\xi)}
 \sum_{n=0}^{\infty}
 \frac{(-2\ii)^n
 (\frac12-\frac{\ii\chi^2}{8})_n
 (-\frac{\ii\chi^2}{8})_n}
 {n!\,|\xi|^{2n}} \notag\\
 &-\frac{\chi B_\chi}{2|\xi|}e^{-\ii\Phi(\xi)}
 \sum_{n=0}^{\infty}
 \frac{(2\ii)^n
 (\frac12+\frac{\ii\chi^2}{8})_n
 (1+\frac{\ii\chi^2}{8})_n}
 {n!\,|\xi|^{2n}}, \tag{26}\\
 \psi_2(\xi)\sim{}&\ \pm B_\chi e^{-\ii\Phi(\xi)}
 \sum_{n=0}^{\infty}
 \frac{(2\ii)^n
 (\frac12+\frac{\ii\chi^2}{8})_n
 (\frac{\ii\chi^2}{8})_n}
 {n!\,|\xi|^{2n}} \notag\\
 &\pm\frac{\chi A_\chi}{2|\xi|}e^{\ii\Phi(\xi)}
 \sum_{n=0}^{\infty}
 \frac{(-2\ii)^n
 (\frac12-\frac{\ii\chi^2}{8})_n
 (1-\frac{\ii\chi^2}{8})_n}
 {n!\,|\xi|^{2n}}. \tag{27}
\end{align}
In particular,
\begin{align}
 \psi_1(\xi)&=A_\chi e^{\ii\Phi(\xi)}
 -\frac{\chi B_\chi}{2|\xi|}e^{-\ii\Phi(\xi)}
 +O(|\xi|^{-2}), \notag\\
 \psi_2(\xi)&=\pm B_\chi e^{-\ii\Phi(\xi)}
 \pm\frac{\chi A_\chi}{2|\xi|}e^{\ii\Phi(\xi)}
 +O(|\xi|^{-2}). \tag{28}
\end{align}
\end{lemma}

\begin{proof}
Apply the standard two-branch expansion of \({}_1F_1\) \cite[Eq.~13.7.2]{DLMF} to (7)--(8), with the principal logarithm of \(\zeta=\ii\xi^2/2\). The exponential and algebraic branches give (25), and the inverse-power factors give the two displayed series.
\end{proof}

The gamma-modulus identities imply
\begin{equation}
 |A_\chi|^2=\frac{1+e^{-\pi\chi^2/4}}2,
 \qquad
 |B_\chi|^2=\frac{1-e^{-\pi\chi^2/4}}2. \tag{29}
\end{equation}
Define the two limiting directions directly by
\begin{equation}
 (u_\pm)^\perp=
 \pm\frac{\chi\pi e^{-\pi\chi^2/8}e^{-\pi\ii/4}}
 {\sqrt2\,\Gamma(\frac12+\frac{\ii\chi^2}{8})
 \Gamma(1-\frac{\ii\chi^2}{8})},
 \qquad
 (u_\pm)_3=e^{-\pi\chi^2/4}. \tag{30}
\end{equation}
The same gamma identities show that \(|u_\pm|=1\). Substitution of (28) into (9)--(10) gives, as \(\xi\to\pm\infty\),
\begin{align}
 \gamma^\perp(\xi)={}&(u_\pm)^\perp
 -\frac{\chi}{\xi}\left(
 A_\chi^2e^{2\ii\Phi(\xi)}
 -B_\chi^2e^{-2\ii\Phi(\xi)}\right)
 +O(|\xi|^{-2}), \notag\\
 \gamma_3(\xi)={}&(u_\pm)_3
 -\frac{2\chi}{|\xi|}\Rea\!\left(
 \overline{A_\chi}B_\chi e^{-2\ii\Phi(\xi)}\right)
 +O(|\xi|^{-2}), \tag{31}\\
 \mathbf d^\perp(\xi)={}&-A_\chi^2e^{2\ii\Phi(\xi)}
 +B_\chi^2e^{-2\ii\Phi(\xi)}
 -\frac{\chi}{\xi}(u_\pm)^\perp+O(|\xi|^{-2}), \notag\\
 (\mathbf d)_3(\xi)={}&\mp2\Rea\!\left(
 \overline{A_\chi}B_\chi e^{-2\ii\Phi(\xi)}\right)
 -\frac{\chi}{\xi}(u_\pm)_3+O(|\xi|^{-2}). \tag{32}
\end{align}
In particular,
\begin{equation}
 \gamma(\xi)\longrightarrow u_\pm,
 \qquad
 \xi\gamma(\xi)-\chi\mathbf d(\xi)=\xi u_\pm+O(|\xi|^{-1}). \tag{33}
\end{equation}

It remains to determine the finite part of \(J_\chi\). Let \(\partial_\chi\) denote differentiation at fixed \(\xi\), and let \(\Psi=\Gamma'/\Gamma\). The Kummer pair (8) satisfies
\begin{equation}
 \psi_1'=\frac{\ii\xi}{2}\psi_1+\frac{\ii\chi}{2}\psi_2,
 \qquad
 \psi_2'=\frac{\ii\chi}{2}\psi_1-\frac{\ii\xi}{2}\psi_2. \tag{34}
\end{equation}
Differentiating (34) with respect to \(\chi\) at fixed \(\xi\), the diagonal terms cancel and give the exact primitive identities
\begin{align}
 J_\chi^\perp&=-2\ii\left(
 \psi_1\partial_\chi\psi_2-
 \psi_2\partial_\chi\psi_1\right), \tag{35}\\
 (J_\chi)_3&=-2\ii\left(
 \overline{\psi_1}\partial_\chi\psi_1+
 \overline{\psi_2}\partial_\chi\psi_2\right), \tag{36}
\end{align}
with \(J_\chi(0)=0\). The logarithmic derivatives of the explicit coefficients (25) are
\begin{align}
 \frac{A_\chi'}{A_\chi}
 &=\frac{\chi}{4}\left[-\frac\pi2-\ii\log2
 -\ii\Psi\!\left(\frac12+\frac{\ii\chi^2}{8}\right)\right], \tag{37}\\
 \frac{B_\chi'}{B_\chi}
 &=\frac1\chi+\frac{\chi}{4}\left[-\frac\pi2+\ii\log2
 +\ii\Psi\!\left(1-\frac{\ii\chi^2}{8}\right)\right]. \tag{38}
\end{align}
Since \(J_\chi'=-\mathbf d\), (32) shows that the nonoscillatory term \(\chi u_\pm/\xi\) produces the logarithm, while \(2\Phi'(\xi)=\xi+\chi^2/(2\xi)\) makes the oscillatory tails \(O(|\xi|^{-1})\) after one integration by parts. Substitution of (28) into (35)--(38) then yields
\begin{equation}
 J_\chi(\xi)=\chi u_\pm\log|\xi|+C_\pm+O(|\xi|^{-1}),
 \qquad \xi\to\pm\infty, \tag{39}
\end{equation}
where the finite parts are entirely explicit:
\begin{align}
 (C_\pm)^\perp={}&\pm
 \frac{\chi\pi e^{-\pi\chi^2/8}e^{-\pi\ii/4}}
 {\sqrt2\,\Gamma(\frac12+\frac{\ii\chi^2}{8})
 \Gamma(1-\frac{\ii\chi^2}{8})} \notag\\
 &\times\left[
 \frac{\ii}{\chi}-\frac{\chi}{4}\left\{
 2\log2+\Psi\!\left(1-\frac{\ii\chi^2}{8}\right)
 +\Psi\!\left(\frac12+\frac{\ii\chi^2}{8}\right)
 \right\}\right], \tag{40}\\
 (C_\pm)_3={}&\frac{\chi}{4}\Bigg[
 \left(1-e^{-\pi\chi^2/4}\right)
 \left\{\log2+\Rea\Psi\!\left(1-\frac{\ii\chi^2}{8}\right)\right\} \notag\\
 &\hspace{18mm}-\left(1+e^{-\pi\chi^2/4}\right)
 \left\{\log2+\Rea\Psi\!\left(\frac12+\frac{\ii\chi^2}{8}\right)\right\}
 \Bigg]. \tag{41}
\end{align}
Thus \((C_-)^\perp=-(C_+)^\perp\) and \((C_-)_3=(C_+)_3\).

\begin{theorem}[Explicit asymptotic lines]
Let the hypotheses of Theorem~1 hold, and set
\begin{equation}
 b_\pm=\rho^{-3/2}\left\{
 cC_\pm-cJ_\chi(\xi_0)
 -a\bigl(\xi_0\gamma(\xi_0)-\chi\mathbf d(\xi_0)\bigr)
 \right\}. \tag{42}
\end{equation}
Then, as \(s\to\pm\infty\),
\begin{align}
 r(s)={}&b_\pm+u_\pm\Bigg[
 \frac{a}{\rho}\left(s+\frac{ab+cd}{\rho^2}\right) \notag\\
 &\hspace{16mm}-\frac{c\Delta}{\rho^3}
 \log\left|\sqrt\rho\left(s+\frac{ab+cd}{\rho^2}\right)\right|
 \Bigg]+O(|s|^{-1}). \tag{43}
\end{align}
Consequently,
\begin{equation}
 \operatorname{dist}\bigl(r(s),\,b_\pm+\mathbb R u_\pm\bigr)
 =O(|s|^{-1}). \tag{44}
\end{equation}
The oriented angle \(\Theta\in[0,\pi]\) between the two asymptotic directions is
\begin{equation}
 \cos\Theta=2e^{-\pi\chi^2/2}-1,
 \qquad
 \Theta=2\arccos\!\left(e^{-\pi\chi^2/4}\right). \tag{45}
\end{equation}
\end{theorem}

\begin{proof}
Insert (33) and (39) into (16), use
\(\xi(s)=\sqrt\rho\bigl(s+(ab+cd)/\rho^2\bigr)\), and recall
\(\chi=-\Delta/\rho^{3/2}\). This gives (43) and the explicit basepoints (42). The logarithmic displacement is parallel to \(u_\pm\), so it does not affect the distance to the corresponding line, proving (44). Finally, (30) and the gamma-modulus identities give
\(u_+\cdot u_-=2e^{-\pi\chi^2/2}-1\), hence (45).
\end{proof}

\paragraph{Geometric interpretation.}
If \(c=0\), the logarithmic term vanishes and (43) is an ordinary affine asymptotic parametrization. If \(a=0\), the longitudinal escape along each asymptotic line is logarithmic. In every nonhelical case \(\Delta\ne0\), the transverse distance to the appropriate line decays as \(O(|s|^{-1})\). The acute angle between the unoriented lines is
\(\arccos|2e^{-\pi\chi^2/2}-1|\).

\paragraph{Numerical illustration.}
Figure~\ref{fig:explicit-example} revisits the example \(\kappa(s)=s+1\), \(\tau(s)=s/10\) used in \cite[Figure~1]{LucasOrtega2026}. Its Cartesian coordinates are evaluated directly from (22)--(23), rather than by numerical integration of the Frenet--Serret equations. The routine \texttt{mpmath.hyper2d} is especially convenient because its parameter dictionaries transcribe (18)--(19) directly; the asymptotic lines are evaluated independently from (30) and (40)--(42).

\begin{figure}[htbp]
\centering
\includegraphics[width=0.78\textwidth]{figure1.pdf}
\caption{The three-dimensional clothoid \(\kappa(s)=s+1\), \(\tau(s)=s/10\), \(-10\le s\le10\), together with its explicit asymptotic lines \(b_-+\mathbb R u_-\) and \(b_++\mathbb R u_+\), shown dashed.}
\label{fig:explicit-example}
\end{figure}

The excluded cases are already explicit: if \(a=c=0\), one obtains the elementary constant-curvature, constant-torsion helix; if \(\Delta=0\) and \((a,c)\ne(0,0)\), then \(\tau/\kappa\) is constant and the commuting clothoid-helix family reduces to Fresnel integrals \cite{Frego2022,Rosu2026}. Thus Theorems~1 and 2 complete both the closed-form Cartesian and asymptotic treatment of affine curvature and torsion.

\clearpage
\appendix
\section{Reproducible evaluation}
The following is the exact Python program used to generate Figure~\ref{fig:explicit-example}.
\lstinputlisting[language=Python]{figure1.py}

\begin{thebibliography}{9}\small
\bibitem{BerryRobbins1993}
M. V. Berry and J. M. Robbins, Classical geometric forces of reaction: an exactly solvable model, \emph{Proc. R. Soc. Lond. A} \textbf{442} (1993), no.~1916, 641--658.
\bibitem{Frego2022}
M. Frego, Closed form parametrisation of 3D clothoids by arclength with both linear varying curvature and torsion, \emph{Appl. Math. Comput.} \textbf{421} (2022), 126907.
\bibitem{Senyurt2015}
S. \c{S}enyurt and A. \c{C}al\i\c{s}kan, Spinor formulation of Sabban frame of curve on \(S^2\), \emph{Pure Math. Sci.} \textbf{4} (2015), no.~1, 37--42.
\bibitem{Ionut2022}
A. Ionu\textcommabelow{t}, On the spherical clothoid, arXiv:2203.07963 [math.DG] (2022).
\bibitem{Jursenas2014}
R. Jur\v{s}\.{e}nas, On the definite integral of two confluent hypergeometric functions related to the Kamp\'e de F\'eriet double series, \emph{Lithuanian Math. J.} \textbf{54} (2014), 61--73.
\bibitem{Rosu2026}
H. C. Rosu, J. de la Cruz, and P. Lemus-Basilio, Clothoid helices obtained via the Lie--Darboux method, arXiv:2604.06577 [math-ph] (2026).
\bibitem{LucasOrtega2026}
P. Lucas and J.-A. Ortega-Yag\"ues, A geometric property characterizing the 3D-generalized clothoids, \emph{Turkish J. Math.} \textbf{50} (2026), no.~2, 240--251.
\bibitem{DLMF}
NIST Digital Library of Mathematical Functions, \S13.7, asymptotic expansions for Kummer functions, release 1.2.7 (2026).
\end{thebibliography}
\end{document}
